\documentclass[12pt,a4paper,fontset=windows]{ctexart}
\usepackage[margin=2.5cm,top=2.55cm,bottom=2.7cm]{geometry}
\usepackage{amsmath}
\setmainfont{Times New Roman}
\linespread{1.18}
\pagestyle{plain}
\begin{document}
\begin{center}\Large\bfseries AI工具使用详情\end{center}
\noindent 本说明记录已知的论文复核与修订过程。工具自动检查与作者人工审查分别说明；尚未收到确认的人工工作不记为完成。

\section*{一、工具名称与用途}
使用OpenAI Codex（GPT-6）辅助阅读用户提供的赛题、解答包、论文源文件和格式参考资料；检查模型、代码和数据的一致性；重绘图表、调整LaTeX排版、编写核验及数据恢复脚本，并生成PDF。求解器使用SciPy调用HiGHS，绘图使用Matplotlib；这些数值工具不替代对模型假设的判断。

\section*{二、主要提示方式与过程}
用户提供第十版论文工作包及解答资料包，要求补充热力图、增加可视化种类、约15个编号图，正文控制在25--29页，并将完整源码移至支撑材料，附录仅保留文件清单与核心算法。随后用户要求再次检查各种格式要求及内容正确性，发现问题后修改。

典型提示包括：“在已生成的论文里面适当补充资料包的热力图，同时不要违反格式要求”；“附录只用列代码文件名，不需要完整源码”；“你确定吗，各种格式要求符合吗，内容正确吗？检查并修改”。AI据此读取已有文件，逐项对照参考资料和官网规范，并按发现的问题修订。后续提示为“认真检查多几轮，不要偷懒”，并要求记住历史复查错误。AI读取所引用任务，提取缺工作表、错时间标签和汇总错误漏报等教训，执行多轮针对性检查。

\newpage
\section*{三、采纳内容与机器核验}
本次沿用已有模型、参数及指定日期主结果，保留15幅编号图，为热力图、预测误差和敏感性分析补充统计区间与解释；将纸面附录整理为文件清单及三个算法摘要。复查补齐本PDF、放大紧急购电表字号、补明单位与预测光伏标记，并使全部图具有可运行重绘入口。

四问主求解入口在隔离副本重跑，五个主情形共192,528条记录与原结果在相应单位的$10^{-5}$绝对容差内一致。问题三的预报误差、更新策略、安全分位数、效率情景及调整快照亦已复算。补充核验覆盖52张数值表、五份结果工作簿、物理约束及费用分项。论文经编译和逐页图像检查；这些是机器复核证据，不表示作者人工审查已经完成，也不证明策略在所有未来情形下全局最优。

新增检查发现：将结果工作簿日初储电量清空后，原核验入口仍通过。修订后必填库存为空或非有限值会明确失败；故障测试同时验证返回码、最终状态和对应错误原因。另以独立变量布局建立代表日线性规划，执行未来信息扰动、24个指定日期复算及费用、库存边界测试。主模型参数与结果工作簿未改动。

\section*{四、人工审查与未完成事项}
用户已给出结构、页数、图数与附录范围的修改意见。尚未收到作者对全部模型假设、原始程序来源、每项AI参与内容和最终论文逐项人工核实完成的确认，因此本说明不声称已完成该环节。此前材料的全部AI使用历史未完整获知，需要作者据实补齐。

管理参数的有限候选坐标搜索未在本轮重跑。退订是否退款、单侧效率解释与模板时间标签对应关系仍按正文所述口径处理。精简附录服从用户要求，但官网第五条要求附录包含完整源码，第十条要求纸电一致；仅将源码放入ZIP不能被本说明认定为豁免。作者应在正式提交前确认这一冲突的处理方式及赛区要求。
\end{document}

